\begin{proposition}[The integral obstruction to a parameter shift]\hypertarget{tag-01Q1}{}\label{fam:src:A.2.1}\label{fam:st:A.2.1}\hypertarget{tag-03CQ}{}\label{wb:pshift:lem:centralizer}
\StatementRecord{fam:src:A.2.1}{proved}{01Q1}{\SeriesRef{fam:src:A.1.2}}

The rational centralizer is
\begin{equation}
 \operatorname{Cent}_{\mathbb Q}(A_1,A_2)
       =\{aI+bE:a,b\in\mathbb Q\}.
 \hypertarget{tag-03BZ}{}\label{wb:pshift:eq:centralizer}
\end{equation}
A biholomorphism over the identity of $B^\circ$ between the torus
families forces $c'-c\in\mathbb Z$ and has complex-linear part
$I_2$ or $-I_2$ on every fibre. Every biholomorphism $X_c\to X_{c'}$ has $c'-c=n\in2\mathbb Z$, complex-linear fibre part $I_2$, and source-to-target lattice map $I-nE$. In the full fixed lattices,
\[
 \Lambda^{A_1}/N_1\Lambda\simeq\mathbb Z/3,\qquad
 \Lambda^{A_2}/N_2\Lambda\simeq\mathbb Z/4\oplus\mathbb Z/2,
\]
with $d$ generating the second factor of the latter group and vanishing in the former.
\end{proposition}
\begin{proof}
\begin{enumerate}[label=\arabic*.,leftmargin=*]
\item For an integer $n$, put $S_n=I+nE$. Since $S_n^{-1}=S_{-n}$,
\begin{equation}
 \Pi_{c+n}=\Pi_c S_n,\qquad
 \Pi_{c+n}\Lambda=\Pi_c\Lambda.
 \hypertarget{tag-03CF}{}\label{wb:pshift:eq:integer-shear}
\end{equation}
Both monodromy matrices fix $\widehat\delta$ and preserve $\gamma$.
Consequently $A_jE=EA_j=E$, and the same commutation holds for every
monodromy word. The identity of $\mathbb C^2$ therefore identifies
the two full middle torus families with their $R_g$-action.
Its source-to-target lattice map is $S_{-n}$, since
$\Pi_{c+n}S_{-n}=\Pi_c$.
\item Solving the two fixed-vector systems from \eqref{wb:pshift:eq:finite-data} gives
\[
 (\Lambda\otimes\mathbb Q)^{A_1}
   =\mathbb Q(1,2,-4,0)^{\mathsf t}+\mathbb Q\widehat\delta,
 \quad
 (\Lambda\otimes\mathbb Q)^{A_2}
   =\mathbb Q(1,3,-3,0)^{\mathsf t}+\mathbb Q\widehat\delta.
\]
Their intersection is $\mathbb Q\widehat\delta$. The fixed covectors
are respectively
\[
 \mathbb Q\gamma+\mathbb Q(2u+w+3\delta),\qquad
 \mathbb Q\gamma+\mathbb Q(u+w+2\delta).
\]
Their intersection is $\mathbb Q\gamma$: comparing the $u$ and $w$
coefficients of a common multiple gives $2x=y$ and $x=y$, hence
$x=y=0$. Every commuting endomorphism therefore preserves
$\mathbb Q\widehat\delta$ and $\ker\gamma$.

\item On $\ker\gamma/\mathbb Q\widehat\delta$, in the classes of
$(\widehat u,\widehat w)$, the two matrices are
\[
 Q_1=\begin{pmatrix}0&1\\-1&-1\end{pmatrix},\qquad
 Q_2=\begin{pmatrix}0&-1\\1&0\end{pmatrix}.
\]
A matrix commuting with $Q_2$ has the form
$\left(\begin{smallmatrix}b&d\\-d&b\end{smallmatrix}\right)$.
Commutation with $Q_1$ gives $d=0$. Thus a commuting endomorphism $H$
has columns
\[
 H\widehat\gamma=(a,r,s,k)^{\mathsf t},\quad
 H\widehat u=b\widehat u+u_0\widehat\delta,\quad
 H\widehat w=b\widehat w+w_0\widehat\delta,\quad
 H\widehat\delta=d_0\widehat\delta.
\]
Use the columns $A_1\widehat u=-\widehat w+\widehat\delta$,
$A_1\widehat w=\widehat u-\widehat w$, and
$A_2\widehat u=\widehat w$. Commutation on these columns gives
\[
 d_0-w_0=b+u_0,\qquad u_0=2w_0,\qquad w_0=u_0.
\]
Hence $u_0=w_0=0$ and $d_0=b$. Commutation on
$\widehat\gamma$ with $A_1$ gives
\[
 r=2(a-b),\qquad s=-4(a-b).
\]
Commutation on that column with $A_2$ also gives $r=-s$.
Therefore $a=b$, $r=s=0$, and $H=aI+kE$. Conversely these matrices
commute, by the fixed-vector and fixed-covector identities preceding
\eqref{wb:pshift:eq:integer-shear}. This proves \eqref{wb:pshift:eq:centralizer}.

\item For the geometric assertion, a holomorphic map between two compact
complex tori lifts to their universal covers. The derivative of the
lift is periodic, so its entries descend to holomorphic functions on
the compact source torus and are constant. The lift is consequently
affine, with a complex-linear part $L$. In a holomorphic family these
linear parts vary holomorphically. Their lattice maps are integral
and locally constant, hence give an isomorphism of the two marked
local systems. Because the map is over the identity of the base,
that integral map $H$ commutes with both $A_j$. Thus
$H=aI+bE$, with $a,b\in\mathbb Z$. Its determinant is $a^4$;
invertibility over $\mathbb Z$ gives $a=1$ or $a=-1$.

\item The complete period equation is
\begin{equation}
                   L(z)\Pi_c(z)=\Pi_{c'}(z)(aI+bE).
 \hypertarget{tag-03CM}{}\label{wb:pshift:eq:period-morphism}
\end{equation}
Its third and fourth columns give $L(z)=aI_2$. Its first column
then gives $a(c'-c)+b=0$. Therefore $c'-c=-b/a\in\mathbb Z$,
and the asserted two linear parts exhaust the possibilities.
\item \emph{The finite-root obstruction.}
The two cyclic norm matrices are, by multiplication of \eqref{wb:pshift:eq:finite-data},
\[
 N_1=I+A_1+A_1^2=
 \begin{pmatrix}3&0&0&0\\6&0&0&0\\-12&0&0&0\\0&2&1&3\end{pmatrix},
\]
\begin{equation}
 N_2=I+A_2+A_2^2+A_2^3=
 \begin{pmatrix}4&0&0&0\\12&0&0&0\\-12&0&0&0\\0&2&2&4\end{pmatrix}.
 \hypertarget{tag-03CK}{}\label{wb:pshift:eq:norms}
\end{equation}
For example these formulas give
$N_1\widehat w=\widehat\delta$,
$N_2\widehat u=N_2\widehat w=2\widehat\delta$,
$N_1\widehat\delta=3\widehat\delta$, and
$N_2\widehat\delta=4\widehat\delta$.

We first justify why a local isomorphism must satisfy the exact affine
equations on the marked finite cover. The ramified base change
$t_j=u^{m_j}$ has local equation
\[
 s_j^{m_j}-u^{m_j}
   =\prod_{a=0}^{m_j-1}(s_j-e^{2\pi ia/m_j}u)=0.
\]
Its normalization separates all these smooth branches, each with
coordinates $(z_1,z_2,u)$. The marked map
$y\mapsto(q_j(y),s_j(y))$ identifies the marked torus family $J'_j$
with this normalization, as proved by the explicit branches in
Proposition~\ref{fam:src:A.1.2}. A biholomorphism over the identity of
$D_j$ base-changes over the identity of the $u$-disc and lifts uniquely
to the normalization: on local rings it carries integral elements in
the total fraction ring to integral elements, and its inverse does
the same. This lift commutes with the marked deck action because
both composites lift the same map $(x,u)\mapsto(F(x),\zeta_j u)$.

Let $c'=c+n$, where $n$ is an integer, and suppose the linear part of
the lifted map is $aI_2$, $a\in\{1,-1\}$. The lifted map on the torus
family is $[z,\zeta]\mapsto[z,a\zeta+B_j(z)]$. Its translating
section has a holomorphic lift $B_j:\Delta_j\to\mathbb C^2$:
$\Delta_j\times\mathbb C^2\to J'_j$ is a covering map and the disc is
simply connected. Local inverses of this covering show that the lift
is holomorphic. The deck-equivariance identity holds modulo a
lattice vector, so, after choosing this lift, it is
\begin{equation}
 B_j(g_jz)-R_{g_j}(z)B_j(z)
  =\Pi_c(g_jz)\left(
       \frac{(1-a)v_j+n\epsilon_j\widehat\delta}{m_j}
          +\ell_j\right),\qquad \ell_j\in\Lambda.
 \hypertarget{tag-03C9}{}\label{wb:pshift:eq:finite-equation}
\end{equation}
The integer column $\ell_j$ is constant: it is locally constant by
discreteness of the marked lattice and the disc is connected.
Iterating $m_j$ times, with $A_jv_j=v_j$ and
$A_j\widehat\delta=\widehat\delta$, yields the necessary equality
\begin{equation}
 N_j\ell_j=-(1-a)v_j-n\epsilon_j\widehat\delta.
 \hypertarget{tag-03CB}{}\label{wb:pshift:eq:finite-norm-obstruction}
\end{equation}
Here the left side of the iterated equation is zero because the
base point and the marked vector coordinate return after $m_j$
steps; real invertibility of $\Pi_c$ makes the lattice equality exact.

If $a=-1$, applying $\gamma$ to \eqref{wb:pshift:eq:finite-norm-obstruction} gives
\[
                   m_j\gamma(\ell_j)=-2\epsilon_j.
\]
This is impossible already at $j=1$, since $3$ does not divide $2$.
Thus a global isomorphism cannot have negative linear part.

For $a=1$, the first row of \eqref{wb:pshift:eq:finite-norm-obstruction} gives $\gamma(\ell_j)=0$.
Equation \eqref{wb:pshift:eq:norms} then proves that the order-three filling imposes no
restriction on $n$, whereas the order-four filling imposes precisely
$n\in2\mathbb Z$. More intrinsically, the finite obstruction is the
class of $-n\epsilon_j\widehat\delta$ in
\begin{equation}
 \Lambda^{A_j}/N_j\Lambda.
 \hypertarget{tag-03C7}{}\label{wb:pshift:eq:finite-class}
\end{equation}
In the marked fixed bases this quotient is
$\mathbb Z/3$ at $j=1$, with $\widehat\delta$ zero, and
$\mathbb Z/4\oplus\mathbb Z/2$ at $j=2$, with
$\widehat\delta$ generating the second summand. Indeed \eqref{wb:pshift:eq:norms} gives
\[
 N_1\Lambda=3\mathbb Z v_1\oplus\mathbb Z\widehat\delta,
 \quad
 N_2\Lambda=4\mathbb Z\varepsilon'\oplus
                      2\mathbb Z\widehat\delta,
 \qquad \varepsilon'=(1,3,-3,0)^{\mathsf t}.
\]
This proves the parity obstruction.
\end{enumerate}
\end{proof}
